\section{Jaweh Adonai}
\begin{block}[Jahweh Adonai]
        Letzes mal haben wir den Namen Jahweh Zebaoth angeschaut. Der Name bedeutet Herr der Heerscharen. Wie wir gelernt haben, zeigt uns der Name, dass Gott nicht nur der Herr über uns Menschen ist, sondern auch über alle Himmeln und die Engel.

        Heute wollen wir ein anderer Name anschauen. Er ist ziemlich speziell und heisst Herr, HERR oder Gott der Herr oder Yahweh Adonai. Der Name kommt ca. 280x in der Bibel vor, davon alleine 200x in Hesekiel. Das erste mal kommt dieser Name in 1. Mose vor. Das letzte Mal in der Offenbarung.

        In 1. Mose 2 11x und 1. Mose 3 7x, wird der Name immer in Verbindung mit uns Menschen genannt.
        \begin{bibelbox}{SCHL}{IMos}{2:4}
            Die ist die Geschichte des Himmels und der Erde, als sie geschaffen wurden, zu der Zeit, als Gott der Herr Erde und Himmel machte.
        \end{bibelbox}
        \begin{bibelbox}{SCHL}{IMos}{2:8}
            Und Gott der Herr pflanzte einen Garten in Eden, im Osten, uns setzte den Menschen dorthin, den er gemacht hatte.
        \end{bibelbox}
        Weil dieser Name nur im Kapitel 2 und am Anfang von Kapitel 3 vor kommt, wurde gesagt, dass es sich um eine andere Schöpfungsgeschichte handelt. Quasi eine zweite Schöpfungsgeschichte von jemand anderem als Moses geschrieben. Wenn wir aber Kapitel 1 und Kapitel 2 anschauen, sehen wir, dass es sich im Kapitel 1 um einen chronologischen Ablauf der Schöpfung handelt und in Kapitel 2 primär um die Schöpfung des Menschen und dann in 3 der Fass des Menschen. Durch die Benutzung des Namens Herr HERR oder Gott der Herr, wird so die Beziehung zwischen dem Menschen und Gott dargestellt.

        Nach dem Kapitel 2 taucht der Name erst wieder in 5 Mose auf wo Mose zu dem Volk Israel spricht:
        \begin{bibelbox}{SCHL}{VMos}{6:4}
            Höre Israel, der Herr ist unser Gott der HERR allein!
        \end{bibelbox}
        Auch hier wieder wird eine Beziehung zwischen Gott und seinem Volk erstellt.
        Anschliessend kommt der bekannte Vers:
        \begin{bibelbox}{SCHL}{VMos}{6:5}
            Und du sollst dein HERRN, deinen Gott, lieben mit deinem ganzen Herzen und mit deiner ganzen Seele und mit deiner ganzen Kraft.
        \end{bibelbox}
        Das ist das Gebot das Jesus wiederholt, als er gefragt wurde, was denn das Höchste Gebot ist. Dort spricht Jesus explizit vom Herrn, dein Gott.
        \begin{bibelbox}{SCHL}{Matt}{22:36-37}
            Meister, welches ist das grösste Gebot im Gesetz? Und Jesus sprach zu ihm: Du sollst den Herrn, deinen Gott, lieben mit deinem ganzen Herzen und mit deiner ganzen Seele und mit deinem ganzen Denken.
        \end{bibelbox}

        Später in Hesekiel dort geht es um das Volk Israel, um seine Wegführung, um die zukünftige Rückführung und das Gericht für das Volk Israel. In diesem Prophetenbuch wird häufig \betonung[b,p]{So spricht Gott der Herr} verwendet. So zum Beispiel:
        \begin{bibelbox}{SCHL}{Hes}{34:15}
            Ich selbst will meine Schafe weiden und sie lagern, spricht Gott, der Herr.
        \end{bibelbox}
        oder
        \begin{bibelbox}{SCHL}{Hes}{36:33}
            So spricht Gott der Herr: Zu jener Zeit, wenn ich euch reinigen werde von allen euren Missetaten, da will ich euch wieder in den Städten wohnen lassen, und die Trümmer sollen wieder aufgebaut werden.
        \end{bibelbox}
        Immer wieder geht es um Gott und den Menschen.

        Im Neuen Testament kommt der Name am Anfang bei der Ankündigung der Geburt Jesus vor
        \begin{bibelbox}{SCHL}{Lk}{1:32}
            Dieser wird gross sein und Sohn des Höchsten genannt werden; und Gott der Herr wird ihm den Thron seines Vaters David geben:
        \end{bibelbox}

        und dann wieder in der Offenbarung
        \begin{bibelbox}{SCHL}{Offb}{22:5}
            Und es wird dort keine Nacht mehr geben, und sie bedürfen nicht eines Leuchters, noch des Lichtes der Sonne, denn Gott, der Herr, erleuchtet sie; uns sie werden herrschen von Ewigkeit zu Ewigkeit.
        \end{bibelbox}
        Auch hier wird Gott der Herr als persönlichen Gott vorgestellt. Es ist Gott mein Herr, der auf der neuen Erde das Licht sein wird.

    \end{block}
    \begin{block}[Ende]

        Dieser Name zeigt die persönliche Beziehung von Gott zu uns Menschen. Gott unser Herr. Der Herr, dein und mein persönlicher Gott. Aber nicht nur Persönlichkeit und Nähe drückt dieser Name aus, sondern auch seine Suveränität, seine Herrschaft, seine Allmacht. Gott ist nicht unser Kumpel. Gott ist unser Herr, dieser Gott gab uns das Leben, dieser Gott schenkt uns das hier alles. Er erschuf Himmel und Erde. 
        
        Der Mensch hat nichts gegen einen Gott, aber als Herr möchte er ihn dann doch nicht haben. Ja, der liebe Gott, der soll schauen das wir glücklich sind und das es uns gut geht. Er soll da sein wenn es mit schlecht geht. Zwischendurch will ich mein Herr sein.

        Die Mehrheit heute sagt, ich bin mein eigener Herr ich brauche keinen Gott. Aber versinkt die Welt nicht immer mehr in ein Chaos ohne, dass dieser Herr und führt und leitet?

        Auch hier in unserer kleinen Gemeinde, ist es wichtig, dass wir uns von diesem unseren Herrn führen lassen. Lasst uns versuchen, gemeinsam seinen Auftrag zu erfüllen wie er im Wort steht
        \begin{bibelbox}{SCHL}{Mk}{16:15}
            Und er sprach zu ihnen: Geht hin in alle Welt und verkündigt das Evangelium der ganzen Schöpfung.
        \end{bibelbox}
        Einfach das Evangelium verkündigen und der Rest überlassen wir dem Herrn.

    \end{block}